\documentclass[../../../include/open-logic-section]{subfiles}

\begin{document}

\olfileid{sth}{card-arithmetic}{opps}
\olsection{定义基本运算}

由于我们无需顾及顺序，基数算术比序数算术更容易定义。
我们将同时定义加法、乘法与指数运算。


\begin{defn}
当 $\cardfont{a}$ 与 $\cardfont{b}$ 是基数时：
\begin{align*}
	\cardfont{a} \cardplus \cardfont{b} &\defis 
	\card{\cardfont{a} \disjointsum \cardfont{b}}\\
	\cardfont{a} \cardtimes \cardfont{b} &\defis 
	\card{\cardfont{a} \times \cardfont{b}}\\
	\cardexpo{\cardfont{a}}{\cardfont{b}} &\defis 
	\card{\funfromto{\cardfont{b}}{\cardfont{a}}}
\end{align*}
其中 $\funfromto{X}{Y} = \Setabs{f}{f \text{ is a function } X \to Y}$。
（容易证明，对任意集合 $X$ 与 $Y$，$\funfromto{X}{Y}$ 都存在；我们将其留作练习。）
\end{defn}

\begin{prob}
在 $\Zminus$ 中证明：对任意集合 $X$ 与~$Y$，$\funfromto{X}{Y}$ 存在。
在 $\ZF$ 中，仿照 \olref[sth][ord-arithmetic][using-addition]{rankcomputation} 从 $\setrank{X}$ 和 $\setrank{Y}$ 计算 $\setrank{\funfromto{X}{Y}}$。
\end{prob}



对此定义稍作解释可能会有所帮助。
关于加法：它使用了不交并 $\disjointsum$（定义见 \olref[ord-arithmetic][add]{defdissum}）；易见，此定义在有穷情形下给出了正确的结果。
关于乘法：\olref[sfr][set][pai]{cardnmprod} 表明，如果 $A$ 有 $n$ 个元素且 $B$ 有 $m$ 个元素，那么 $A \times B$ 有 $n \cdot m$ 个元素，因此我们的定义无非是将该思想推广到超穷乘法。
指数运算亦然：我们只是将有穷情形的思想推广到超穷情形。
事实上，在某些方面，超穷基数算术看起来比超穷序数算术更像是“普通”算术：


\begin{prop}\ollabel{cardplustimescommute}
$\cardplus$ 与 $\cardtimes$ 满足交换律与结合律。
\end{prop}

\begin{proof}
对于交换律，根据 \olref[cardinals][cardsasords]{lem:CardinalsBehaveRight}，只需观察到 $\cardeq{(\cardfont{a} \disjointsum
\cardfont{b})}{(\cardfont{b} \disjointsum \cardfont{a})}$ 且 $\cardeq{(\cardfont{a} \times \cardfont{b})}{(\cardfont{b} \times
\cardfont{a})}$。
我们将结合律的证明留作练习。
\end{proof}

\begin{prob}
证明 $\cardplus$ 与 $\cardtimes$ 满足结合律。
\end{prob}

\begin{prop}
$A$ 是无穷的当且仅当 $\card{A} \cardplus 1 = 1 \cardplus \card{A} = \card{A}$。
\end{prop}

\begin{proof}
同 \olref[cardinals][classing]{generalinfinitycharacter} 的证明，由 \olref[ord-arithmetic][using-addition]{ordinfinitycharacter} 与 \olref[cardinals][cardsasords]{lem:CardinalsBehaveRight} 即得。
\end{proof}



这就解释了为何序数的加法/乘法与基数的加法/乘法需要使用不同的符号：它们是\emph{真正不同}的运算。
接下来的两个结果表明，序数的指数运算与基数的指数运算也是不同的运算。
（回顾 \olref[z][infinity-again]{defnomega} 可知 $2 = \{0, 1\}$）：


\begin{lem}\ollabel{lem:SizePowerset2Exp}
$\card{\Pow{A}} = \cardexpo{2}{\card{A}}$，对任意集合 $A$ 成立。
\end{lem}

\begin{proof}
对每个子集 $B \subseteq A$，令 $\chi_B \in \funfromto{A}{2}$ 由下式给出：
\begin{align*}
	\chi_{B}(x) &\defis
	\begin{cases}
		1 & \text{若 }x\in B\\
		0 & \text{否则。}
	\end{cases}
\end{align*}
现在令 $f(B) = \chi_B$；这便定义了一个双射 $f \colon \Pow{A}
\to \funfromto{A}{2}$。所以 $\cardeq{\Pow{A}}{\funfromto{A}{2}}$。因此
$\cardeq{\Pow{A}}{\funfromto{\card{A}}{2}}$，于是
$\card{\Pow{A}} = \card{\funfromto{\card{A}}{2}} =
2^{\card{A}}$。
\end{proof}



这个利落的证明本质上涵盖了 \olref[sfr][siz][red-alt]{sec} 中的讨论。
在那里，我们展示了如何将 $\Pow{\omega}$ 的不可数性“归约”为无穷二进制字符串集 $\Bin^\omega$ 的不可数性。
实际上，$\Bin^{\omega}$ 就是 $\funfromto{\omega}{2}$；且上述证明表明，我们在 \olref[sfr][siz][red-alt]{sec} 中进行的推理，对于用任意集合~$A$ 替代~$\omega$ 的情形同样成立。
该结果也导出了关于基数指数运算的一个简明推论：


\begin{cor}\ollabel{cantorcor}
$\cardfont{a} < \cardexpo{2}{\cardfont{a}}$ 对任意基数~$\cardfont{a}$ 成立。
\end{cor}

\begin{proof}
由康托尔定理（\olref[sfr][siz][car]{thm:cantor}）与 \olref{lem:SizePowerset2Exp} 即得。
\end{proof}



所以 $\omega < \cardexpo{2}{\omega}$。但请注意：这是关于\emph{基数}指数运算的结果。它应当与\emph{序数}指数运算对比，因为在后一种情形下 $\omega = \ordexpo{2}{\omega}$（参见 \olref[ord-arithmetic][expo]{sec}）。

借此讨论基数指数运算的机会，我们也可以对 $\Real$ 的“不可数方式”作出更精确的刻画。


\begin{thm}\ollabel{continuumis2aleph0}
$\card{\Real} = \cardexpo{2}{\omega}$
\end{thm}

\begin{proof}[证明概要]
证明方法有很多。最直接的方法是论证 $\cardle{\Pow{\omega}}{\Real}$ 且 $\cardle{\Real}{\Pow{\omega}}$，然后利用施罗德-伯恩斯坦定理推断 $\cardeq{\Real}{\Pow{\omega}}$，并借助 \olref[card-arithmetic][opps]{lem:SizePowerset2Exp} 推断 $\card{\Real} = \cardexpo{2}{\omega}$。
我们将构造单射 $f \colon \Pow{\omega} \to \Real$ 和 $g \colon \Real \to \Pow{\omega}$ 留作（富有启发性的）练习。
\end{proof}

\begin{prob}
通过证明 $\cardle{\Pow{\omega}}{\Real}$ 与 $\cardle{\Real}{\Pow{\omega}}$，完成 \olref[sth][card-arithmetic][opps]{continuumis2aleph0} 的证明。
\end{prob}

\end{document}